%Runic poetry from Denmark and Scania.

\section{Introduction}

The signa used for Danish inscriptions vary.  According to the more common system all Danish inscriptions are numbered according to a unitary system based on the Danish corpus edition, which unlike the Swedish one is not divided into provinces.  The online database \emph{Danske Runeindskrifter} (‘Danish Runic Inscriptions’) does however follow a system identical to the Swedish one, and since those signa are easier to work with due to the ability to immediately identify the origin of any given inscription it has been adopted here, although references to the DR signa are given in parentheses following the origin.

\subsection{Abbreviations of Danish regions}

\begin{itemize}
  \item Bh = Bornholm
  \item MJy = Midtjylland (Middle Jutland)
  \item NJy = Nordjylland (North Jutland)
  \item Sj = Sjælland (Zealand)
  \item SJy = Sønderjylland (South Jutland)
  \item SlB = Slesvig By (the town of Slesvig)
  \item Sk = Skåne (Scania)
  \item Syd = Lolland-Falster

\end{itemize}

\sectionline

\newpage


\section{Bh 33 (DR 379; Nylarsker-sten 1)}

\begin{flushright}%
\textbf{Dating:} C11th–12th

\textbf{Meter:} \Fornyrdislag
\end{flushright}%

\subsection{Introduction}

TODO.

\subsection{Transliteration}

\textbf{sasur · lit · resa · sten · eftiʀ · aluarþ · faþur · sin · truknaþi · han · uti · meþ · ala · sk¶ibara · etki · i · kristr · hab¶(i) ¶ siolu ¶ has ¶ sten · þesi · stai ¶ eftir}

\subsection{Text}

\bpg\bpa[0]%
\startsub Sassurr lét rêsa stên ęptiʀ Hall·varð, fǫður sinn.\endsub \\
Druknaði hann \alst{ú}ti \hld\ með \alst{a}lla skipara. \\
\textbf{†etki†} é Kristr hialpi siólu hans. \\
\alst{St}ênn þęssi \hld\ \alst{st}ái ęptiʀ.\epa

\bpb Saswer had this stone raised after Hallward, his father. \\
He drowned far out with all the shipmen. \\
\dots may Christ always help his soul. \\
May this stone afterwards stand.\epb\epg

\sectionline\newpage


\section{Bh 34 (DR 380; Nylarsker-sten 2)}

\begin{flushright}%
\textbf{Dating:} C11th–12th

\textbf{Meter:} \Fornyrdislag
\end{flushright}%

\subsection{Introduction}

TODO.

\subsection{Transliteration}

\textbf{kobu:suain : raisti : stain : þ(e)na : a(f)tir : bausa : sun : sin : tr(i)... ...n : þan : is : tribin : ua(r)þ : i : (u)(r)ostu : at : ut:la(n)(k)iu : kuþ : tr(u)tin : hi(a)lbi : hans : ont : auk : sata : mikial :}

\subsection{Text}

\bpg\bpa[0]%
\startsub Kǫ́pu-Svęinn ręisti stęin þęnna ęptir Bausa, sun sinn,\endsub \\
\alst{d}ręng góðan, \hld\ þann es \alst{d}repinn varð \\
ï \alst{o}rrostu \hld\ at \alst{Ú}t-lęngju. \\
Guð dróttinn hialpi hans ǫnd auk \edtrans{santa Mikael}{Saint Michael}{\Bfootnote{Something of a patron saint on Bornholm, also invoked in Bh 1, Bh 4, and Bh 49.  Cf. Syd 11 below.}}.\epa

\bpb Cloak-Swain raised this stone after Beaze, his son, \\
the good dreng who was slain \\
in battle by Utlengie. \\
Lord God and Saint Michael help his spirit.\epb\epg


\section{Bh 42 (DR 387; Vestermarie-sten 5)}

\begin{flushright}%
\textbf{Dating:} C11th–12th

\textbf{Meter:} \Ljodahattr
\end{flushright}%

\subsection{Introduction}

TODO.

\subsection{Transliteration}

\textbf{: asualdi : risti : stein : þinsa : iftʀ : alfar : bruþur : sin : drinr : koþr ¶ : trebin u:syni : auk : skogi : suek : saklausan :}

\subsection{Text}

\bpg\bpa[0]%
\startsub Ǫ̇s·valdi rísti stęin þęnnsa ęptiʀ Alf·ar, bróður sinn,\endsub \\
\alst{d}ręngr góðr \hld\ \alst{d}repinn u̇·sẏni, \\
\ind auk Skógi \alst{s}vęik \alst{s}ak-lausan.\epa

\bpb Oswald carved this stone after Elver, his brother; \\
the good dreng [was] shamefully slain \\
\ind and Skoog betrayed him, guiltless.\epb\epg

\sectionline\newpage


\section{MJy 37 (DR 94; Ålum-sten 1)}

\begin{flushright}%
\textbf{Dating:} C11th

\textbf{Meter:} \Fornyrdislag
\end{flushright}%

\subsection{Introduction}

TODO.

\subsection{Transliteration}

\textbf{tuli : (r)(i)s-(i) : stin : þasi : aft ¶ ikal:t : sun : sin : miuk (:) (k)¶(u)... ...k : þau : mun(u) ¶ mini : m-(r)gt : ¶ iuf (:) þirta :}

\subsection{Text}

\bpg\bpa[0]%
\startsub Tóli rêsþi stên þannsi apt Ingj·ald sun sinn, miok gó\emph{ðan dręn}g.\endsub \\
Þau \alst{m}unu \alst{m}inni \hld\ \alst{m}\emph{ę}rkt ê of \emph{b}irta.\epa

\bpb Toole raised this stone after Ingeld, his son, a very good dreng. \\
TODO.\epb\epg


\sectionline\newpage

\section{MJy 71 (DR 131; Års-sten)}

\begin{flushright}%
\textbf{Dating:} C11th

\textbf{Meter:} \Fornyrdislag
\end{flushright}%

\subsection{Introduction}

TODO.

\subsection{Transliteration}

§ A \textbf{: ąsur : sati : stin : þąnsi : aft : ual:tuka : trutin : ¶ × sin} \\
§ B \textbf{× stin : kuask : hirsi : stąnta : ląki : saʀ : ual:tuka × ¶ × uarþa : nafni}

\subsection{Text}

\bpg\bpa[0]%
Ǫnsurr satti stên þannsi apt Val-Tóka, dróttin sinn. \\
\alst{St}ênn kweðsk hérsi \hld\ \alst{st}anda lęngi; \\
sáʀ \alst{V}al-Tóka \hld\ \edtrans{\alst{v}arða}{cairn}{\Bfootnote{Acc. sg. of \emph{varði} ‘cairn or beacon raised as a landmark’.  The word still survives in Danish \emph{varde} ‘id.’, Faroese \emph{varði} ‘id.’}} nęfni.\epa

\bpb Answer set this stone after Wal-Took, his lord. \\
The stone proclaims it will stand here for long; \\
may it name Wal-Took’s cairn.\epb\epg

\sectionline\newpage


\section{MJy 79 (DR 68; Århus-sten 5)}

\begin{flushright}%
\textbf{Dating:} C10th–11th

\textbf{Meter:} \Ljodahattr
\end{flushright}%

\subsection{Introduction}

A well-preserved stone with a small piece struck off at the lower corner at the very beginning of the inscription.

\subsection{Transliteration}

§ A \textbf{(-)usti × auk × hufi × auk × þir × frebiurn × risþu × stin × þąnsi × eftiʀ × ¶ × ąsur × saksa × filaka × sin × harþa ×} \\
§ B \textbf{kuþan × trik × saʀ × tu × ¶ × mana × mest × uniþikʀ × ¶ saʀ × ati × skib × miþ × arną +}

\subsection{Text}

\bpg\bpa[0]%
\startsub \emph{T}osti auk Hófi þêʀ Frø̂·biorn rêsþu stên þannsi ęptiʀ Ǫnsur saxa, fé-laga sinn, harða góðan dręng.\endsub \\
Sáʀ dó \alst{m}anna \hld\ \alst{m}ęst u̇·níðingʀ; \\
\ind sáʀ \alst{á}tti skip með \alst{A}rna.\epa

\bpb Toste and Hove, they and Freebern raised this stone after Answer the seax, his partner, a very good dreng. \\
He died the greatest un-nithing of men; \\
\ind he owned a ship with Arnie.\epb\epg

\sectionline\newpage


\section{NJy 41 (DR 149; Ydby-sten)}

\begin{flushright}%
\textbf{Dating:} C10th–11th

\textbf{Meter:} \Fornyrdislag
\end{flushright}%

\subsection{Introduction}

The inscription is entirely poetic.  The stone is heavily damaged and divided into several parts but the inscription survives in older depictions.

\subsection{Transliteration}

§ A \textbf{þurkisl : [sati · au(k) (·) suniʀ (·) lifa · i] : s[t]aþ [· þa]¶:nsi :} \\
§ B \textbf{[sti](n) (÷) (u)[ftiʀ · lifa]}

\subsection{Text}

\bvg\bva[]%
Þȯr·gísl \alst{s}atti \hld\ ok \alst{s}yniʀ Lêfa &
ï \alst{st}að þannsi \hld\ \alst{st}ên øptiʀ Lêfa. \eva

\bvb Thuryils set and the sons of Lofe \\
in this place the stone after Lofe.\evb\evg

\sectionline\newpage


\section{Sj 93 (DR 229; Sandby-sten 3)}

\begin{flushright}%
\textbf{Dating:} C11th

\textbf{Meter:} \Fornyrdislag
\end{flushright}%

\subsection{Introduction}

A heavily damaged stone, the verse can be tentatively reconstructed based on Syd 11.

\subsection{Transliteration}

§ A \textbf{sylfa : rest... ... ...i : sbalklusu : eyfti : susur : faþur ... ...(r)þi : bru : þisi : iki : ¶ þurils : brþur : sin :} \\
§ B \textbf{i mun · san... ¶ ...if · uitrik · susi · eʀ · uan · sil...}

\subsection{Text}

\bpg\bpa[0]%
\startsub Sylfa rêst\emph{i stên þannsi} ï Spalk-lø̂su øpti Susur, fǫður \emph{sinn, ok gę}rði bró þęssi ę\emph{pt}i Þȯr·ils, br\emph{ó}ður sinn.\endsub \\
Ê mun \alst{s\emph{t}}an\emph{da, \hld\ með \alst{st}ênn á l}íf, \\
\alst{v}itring súsi, \hld\ eʀ \alst{v}ann Syl\emph{fa}. \epa

\bpb Silve raised this stone in Spalklease after Saswer, his father, and made this bridge after Thuryils, his brother. \\
Ever will stand while the stone has life \\
this memento which Silve accomplished.\epb\epg

\sectionline\newpage


\section{SJy 5 (DR 37; Egtved-sten)}

\begin{flushright}%
\textbf{Dating:} C10th–11th

\textbf{Meter:} Unclear
\end{flushright}%

\subsection{Introduction}

A fragmented stone; the remaining piece is easily read, but a large part is missing.

\subsection{Transliteration}

\textbf{... ... ...at ' fai(n) ['] (t)u ÷ i suiu ' raist ¶ ... ...uþiʀ ' aft ' bruþur ¶ stain ' sasi ' skarni ' ...}

\subsection{Text}

\bvg\bva[]%
\dots\ at fáinn; dó ï Svíu. &
\alst{R}ęist \emph{iak \alst{r}u̇naʀ,} &
\emph{\alst{b}r}óðiʀ apt \alst{b}róður; &
stęin sási skęrni \dots\eva

\bvb the painted; died in Sweden. \\
I carved the runes, \\
a brother after [my] brother; \\
may this stone TODO \dots\evb\evg

\sectionline\newpage


\section{SJy 14 (DR 40; Randbøl-sten)}

\begin{flushright}%
\textbf{Dating:} C10th

\textbf{Meter:} \Ljodahattr
\end{flushright}%

\subsection{Introduction}

A rare instance of a stone carved by a man for his wife.

\subsection{Transliteration}

\textbf{tufi ÷ bruti ÷ risþi ÷ stin ÷ þansi ÷ aft ÷ lika ÷ ¶ brutia ÷ þiʀ ÷ stafaʀ ÷ munu ÷ ¶ þurkuni ÷ miuk ÷ liki ÷ lifa ÷}

\subsection{Text}

\bpg\bpa[0]%
\startsub Tófi bryti rêsþi stên þannsi apt \edtrans{líka}{helpmate}{\Bfootnote{Or ‘equal, match’.  Although grammatically masculine (the nom. sg. is \emph{líki}) the word can be used for both genders, as here.}} brytja.\endsub \\
\alst{Þ}êʀ stafaʀ \hld\ munu \alst{Þ}ȯr·gunni \\
\ind miok \alst{l}ęngi \edtrans{\alst{l}ifa}{live}{\Bfootnote{Stones and staves ‘living’ is a common expression.}}.\epa

\bpb Tove the bailiff raised this stone after the bailiff’s helpmate. \\
These staves will for Thurguth \\
\ind for a very long time live.\epb\epg

\sectionline\newpage


\section{Sk 5 (DR 263; Skabersjö-fibula)}

\begin{flushright}%
\textbf{Dating:} C8th

\textbf{Meter:} \Fornyrdislag
\end{flushright}%

\subsection{Introduction}

TODO.

\subsection{Transliteration}

§ P \textbf{    ... ʀʀʀʀʀʀʀʀʀʀʀʀʀʀʀʀ -aþi tuk fauka fiaʀ sis in a iak as(u) (þ)ui launat ¶ ... auab-k sua fakat...} \\
§ Q \textbf{    ... (ʀ)(ʀ)ʀʀʀʀʀʀʀʀʀʀʀʀʀʀ (r)aþi (t)(u)k fauka fiaʀ sis in a iak asa þui launat ¶ ... auab-k sua fakat...}

\subsection{Text}

\bpg\bpa[0]%
TODO: text\epa

\bpb TODO: trans\epb\epg

\sectionline\newpage


\section{Sk 42 (DR 337; Valleberga-sten)}

\begin{flushright}%
\textbf{Dating:} C11th

\textbf{Meter:} \Fornyrdislag
\end{flushright}%

\subsection{Introduction}

A stone.  TODO.

\subsection{Transliteration}

§ A \textbf{: suin : auk : þurgutr : kiaurþu : kubl : þisi ¶ eftiʀ : mana ¶ auk · suina} \\
§ B \textbf{kuþ : hialbi : siaul : þeʀa : uel : ian : þeʀ : likia : i : luntunum}

\subsection{Text}

\bpg\bpa[0]%
\startsub Svênn auk Þȯr·gø̂tr giǫrðu kumbl þęssi ęptiʀ Manna auk Svenna. Guð hialpi siǫ́l þêʀa vęl,\endsub \\
ian þêʀ \alst{l}iggja \hld\ ï \alst{L}ondúnum.\epa

\bpb Swayn and Thurgeat made these monuments after Manny and Swenny.  God help their soul well, \\
but they lie in London.\epb\epg

\sectionline\newpage


\section{Sk 51 (DR 279; Sjörup-sten)}

\begin{flushright}%
\textbf{Dating:} C11th

\textbf{Meter:} \Ljodahattr
\end{flushright}%

\subsection{Introduction}

TODO.

The incorrect use of \textbf{h} in \textbf{huptiʀ} \emph{øptiʀ} and the lack of it in \textbf{an} \emph{’ann} (= ON \emph{hann} ‘he’) \textbf{a(f)þi} \emph{’afði} (= ON \emph{hafði} ‘had’) indicates the common East Norse loss of that phoneme.

\subsection{Transliteration}

\textbf{[+ sa]ksi : sati : st[in] : þasi : huftiʀ : ą[s]biurn : (s)in : fil(a)gą ' ¶ (t)u-a[s : sun :] ¶ saʀ : flu : aki : a[t :] ub:sal(u)m : an : ua : maþ : an : u¶abn : a(f)þi '}

\subsection{Text}

\bpg\bpa[0]%
\startsub Saxi satti stên þa\emph{nn}si øptiʀ Ǫ̇s·biorn, sinn fé-laga, Tó\emph{k}a* sun.\endsub \\
Sáʀ fló \alst{ê}gi \hld\ at \alst{U}pp-sǫlum, \\
\ind ęn \alst{v}á, með ’ann \alst{v}ǫ́pn ’afði.\epa

\bpb Saxe set this stone after Osbern, his partner, Took’s son. \\
He did not flee at Upsele \\
\ind but fought as long as he had weapons.\epb\epg

\sectionline\newpage


\section{Sk 80 (DR 295; Hällestad-sten 1)}

\begin{flushright}%
\textbf{Dating:} C11th

\textbf{Meter:} \Fornyrdislag
\end{flushright}%

\subsection{Introduction}

TODO.

\subsection{Transliteration}

§ A \textbf{: askil : sati : stin : þansi : ift[iʀ] ¶ : tuka : kurms : sun : saʀ : hulan : ¶ trutin : saʀ : flu : aigi : at : ub:¶:salum} \\
§ B \textbf{satu : trikaʀ : iftiʀ : sin : bruþr ¶ stin : ą : biarki : stuþan : runum : þiʀ :} \\
§ C \textbf{(k)(u)(r)(m)(s) (:) (t)(u)(k)(a) : kiku : (n)(i)(s)(t)[iʀ]}

\subsection{Text}

\bpg\bpa[0]%
\startsub Ǫ̇s·kęll satti stên þannsi ęptiʀ Tóka, Gorms sun, séʀ hollan dróttin;\endsub \\
sáʀ fló \alst{ęi}gi \hld\ at \alst{U}pp-sǫlum. \\
\alst{S}ǫttu dręngaʀ \hld\ ęptiʀ \alst{s}inn bróður \\
\alst{st}ên ȧ biargi \hld\ \alst{st}uððan ru̇num; \\
þêʀ \alst{G}orms Tóka \hld\ \alst{g}ingu nę́stiʀ. \epa

\bpb Oskettle set this stone after Took, Gorm’s son, to him a trusty lord; \\
he did not flee at Upsele. \\
The drengs set after their brother \\
a stone on the hill propped up by runes; \\
they marched closest to Gorm’s Took.\epb\epg

\sectionline\newpage


\section{Sk 81 (DR 296; Hällestad-sten 2)}

\begin{flushright}%
\textbf{Dating:} C11th

\textbf{Meter:} \Fornyrdislag
\end{flushright}%

\subsection{Introduction}

TODO.

\subsection{Transliteration}

\textbf{: ąskautr : ristþi : stin : þansi (:) ¶ (:) (i)ftiʀ : airu : brþur : sin : ian : ¶ : saʀ : uas : him:þiki : tuka : nu : ¶ : skal : statą : stin : ą : biarki :}

\subsection{Text}

\bpg\bpa[0]%
\startsub Ǫ̇s·gautr rêsþi stên þannsi ęptiʀ Ęiru, bróður sinn, ian sáʀ vas hêm-þęgi Tóka.\endsub
Nú skal \alst{st}anda \hld\ \alst{st}ên ȧ biargi.\epa

\bpb Osgeat raised this stone after Ore, his brother, and he was Took’s retainer. \\
Now shall stand the stone on the hill.\epb\epg

\sectionline\newpage


\section{SlB 3 (DR Schl3; Slesvig-pind)}

\begin{flushright}%
\textbf{Dating:} C11th

\textbf{Meter:} \Ljodahattr
\end{flushright}%

\subsection{Introduction}

A rune-inscribed stick containing a puerile curse.

\subsection{Transliteration}

§ A \textbf{runaʀ ÷ iag ÷ risti ÷ a ÷ rikiata ÷ tre sua} \\
§ B \textbf{reþ ÷ saʀ ÷ riki ÷ mo(g)ʀ ÷ asiʀ ÷ a ÷ artagum} \\
§ C \textbf{hulaʀ ÷ auk ÷ bulaʀ ÷ meli ÷ þeʀ} \\
§ D \textbf{ars ÷ sum ÷ magi}

\subsection{Text}

\bvg\bva[]%
\alst{R}u̇naʀ iak \alst{r}ísti \hld\ ȧ \alst{r}íkjanda tré &
\ind svá \alst{r}éð sáʀ \alst{r}íki mǫgr; &
\alst{ę̇}siʀ ȧ \alst{á}r-dǫgum \hld\ hullaʀ ok bullaʀ &
\ind \alst{m}ę́li þéʀ ars sum \alst{m}agi.\eva

\bvb Runes I carve on the powerful wood; \\
\ind so did the powerful boy counsel: \\
Eese in days of yore, hurlies and burlies, \\
\ind may thy arse speak like thy belly!\evb\evg

\sectionline\newpage


\section{Syd 11 (DR 212; Tillitse-sten)}

\begin{flushright}%
\textbf{Dating:} C11th

\textbf{Meter:} \Fornyrdislag
\end{flushright}%

\subsection{Introduction}

The stone actually has three sides, but § C is by a different carver and dedicated to a different person; as such it ought to be treated as a separate inscription and not included.

\subsection{Transliteration}

§ A \textbf{eskil : sulka : sun : let : res(a) ¶ sten : þena : eft : sialfan ¶ sik · e mun · stanta · meþ · sten ¶ lifiʀ · uitrint · su · iaʀ · uan · eskil} \\
§ B \textbf{kristr · hialbi · siol · hans ¶ aok · santa · migael}

\subsection{Text}

\bpg\bpa[0]%
\startsub Ę̇s·kil, Súlka sun, lét rêsa stên þęnna ępt sialfan sik.\endsub \\
É mun \alst{st}anda, \hld\ með \alst{st}ênn lifiʀ, \\
\alst{v}itrind sú \hld\ iaʀ \alst{v}ann Ę̇s·kil. \\
Kristr hialpi siól hans auk santa Mikael.\epa

\bpb Oskettle, Sulk’s son, had this stone raised after himself. \\
Ever will stand while the stone lives \\
this memento which Oskettle accomplished. \\
Christ and Saint Michael help his soul.\epb\epg

\sectionline
